\subsection{INVARIANT BILINEAR FORMS}

PROPOSITION 5. Let \((\mathfrak{g},\mathfrak{h})\) be a split semi-simple Lie algebra, \(\Phi\) an invariant symmetric bilinear form on g, and W the Weyl group of \((\mathfrak{g},\mathfrak{h})\) . Then the restriction \(\Phi'\) of \(\Phi\) to h is invariant under W. Moreover, if \(\Phi\) is nondegenerate, so is \(\Phi'\) .

Let \(\alpha \in \mathbb{R}\) , let \(X_{\alpha}\) be a non-zero element of \(\mathfrak{g}^{\alpha}\) , \(\rho\) the associated representation of \(\mathfrak{sl}(2,k)\) on \(\mathfrak{g}\) , and \(\pi\) the representation of \(\mathbf{SL}(2,k)\) on \(\mathfrak{g}\) compatible with \(\rho\) . Then \(\varPhi\) is invariant under \(\rho\) , and hence under \(\pi\) (§1, no. 4). In particular, \(\varPhi'\) is invariant under \(\theta_{\alpha}(t)|\mathfrak{h}\) (no. 2), and hence under W. The last assertion follows from Prop. 1 (i).

PROPOSITION 6. Let \((\mathfrak{g},\mathfrak{h})\) be a split semi-simple Lie algebra, \(\varPhi\) a nondegenerate invariant symmetric bilinear form on \(\mathfrak{g}\) . For all \(\alpha \in \mathbb{R}\) , let \(X_{\alpha}\) be a non-zero element of \(\mathfrak{g}^{\alpha}\) . Let \((H_i)_{i\in I}\) be a basis of \(\mathfrak{h}\) , and \((H_i')_{i\in I}\) the basis of \(\mathfrak{h}\) such that \(\varPhi(H_i,H_j') = \delta_{ij}\) . The Casimir element associated to \(\varPhi\) in the enveloping algebra of \(\mathfrak{g}\) (Chap. I, §3, no. 7) is then

\[
\sum_ {\alpha \in \mathrm{R}} \frac {1}{\varPhi (X _ {\alpha} , X _ {- \alpha})} X _ {\alpha} X _ {- \alpha} + \sum_ {i \in I} H _ {i} H _ {i} ^ {\prime}.
\]

Indeed, by Prop. 1, \(\Phi(H_i, X_\alpha) = \Phi(H_i', X_\alpha) = 0\) for all \(i \in I\) , \(\alpha \in \mathrm{R}\) , and \(\Phi\left(\frac{1}{\Phi(X_\alpha, X_{-\alpha})} X_\alpha, X_{-\beta}\right) = \delta_{\alpha\beta}\) for all \(\alpha, \beta \in \mathrm{R}\) .

